\documentclass[aspectratio=169,14pt]{beamer}
\usepackage{common}


\title{Программирование}
\subtitle{Организационные вопросы}
\author{a.glushko@g.nsu.ru}
%\institute{University of ShareLaTeX}
\date{6 февраля 2021}


\begin{document}

    \begin{frame}
        \titlepage
    \end{frame}

    \begin{frame}{Преподаватель}
        \begin{block}{Глушко Алексей Вячеславович}
            \begin{itemize}
                \item a.glushko@g.nsu.ru
                \item \includegraphics[height=\fontcharht\font`\B]{media/telegram128} \url{https://t.me/haitaka} % TODO icon
                \item +7~(996)~381~38~04
            \end{itemize}
        \end{block}
    \end{frame}

    \begin{frame}{Группа в Telegram}
        %\begin{center}
        \includegraphics[width=4cm]{media/telegram-invite-qr}

        \includegraphics[height=1em]{media/telegram128} \url{https://t.me/joinchat/SKNwjlPIOPBPn6zK} % TODO icon
        %\end{center}
    \end{frame}

    \begin{frame}{Курс}
        \begin{itemize}
            \item Язык C
            \item Алгоритмы и структуры данных
                \begin{itemize}
                    \item На языке C
                \end{itemize}
        \end{itemize}
    \end{frame}

    \begin{frame}{Условия зачёта}
        \begin{columns}
            \column{0.43\textwidth}
                \begin{itemize}
                    \item максимум 100 баллов
                    \item $\geqslant 50$ --- оценка ``удовл''
                    \item $\geqslant 75$ --- оценка ``хорошо''
                    \item $\geqslant 85$ --- оценка ``отлично''
                \end{itemize}
            \column{0.6\textwidth}
                \begin{itemize}
                    \item Теория: 50 баллов
                    \item Практика: 50 баллов
                        \begin{itemize}
                            \item $5~\text{задач} \times 7~\text{баллов} = 35~\text{баллов}$
                            \item $2~\text{контрольных} \times 5~\text{баллов} = 10~\text{баллов}$
                            \item $5$ баллов дополнительно
                        \end{itemize}
                \end{itemize}
        \end{columns}
    \end{frame}

    \begin{frame}{Задачи}
        \begin{itemize}
            \item 5 задач
            \item 3 уровня сложности
                \begin{itemize}
                    \item сложная --- 6 баллов
                    \item средняя --- 5 баллов
                    \item лёгкая --- 4 балла
                \end{itemize}
            \item несколько вариантов каждого уровня
            \item Система тестирования \\
                \url{https://olympic.nsu.ru/nsuts-new}
        \end{itemize}
    \end{frame}

    \begin{frame}{Дедлайны}
        \begin{itemize}
            \item 2 недели на прохождение автоматических тестов
            \item Ещё 2 недели на защиту (на семинаре?)
            \item Награды
                \begin{itemize}
                    \item Защита в первые две недели: $+1$ балл
                \end{itemize}
            \item Штрафы
                \begin{itemize}
                    \item Прохождение тестов в третью неделю: $-1$ балл
                    \item Прохождение тестов в четвёртую неделю: $-2$ балла
                    \item После 4 недель задача не принимается
                \end{itemize}
        \end{itemize}
    \end{frame}

    \begin{frame}{Досрочный экзамен}
        \begin{itemize}
            \item Принимает семинарист
            \item 5 сложных задач
            \item Вовремя
            \item Со всеми доп. заданиями % TODO
        \end{itemize}
    \end{frame}

    \begin{frame}{Опрос}
        \begin{itemize}
            \item Опыт программирования
                \begin{itemize}
                    \item Школа? Дополнительно?
                    \item Языки
                    \item Участие в олимпиадах
                \end{itemize}
            \item Как в 8 битах представляется число $-1$?
            \item Какие алгоритмы сортировки вы знаете?
        \end{itemize}
    \end{frame}

    \begin{frame}{Демонстрация}
        \begin{itemize}
            \item Системы тестирования
            \item IDE
        \end{itemize}
    \end{frame}

\end{document}
